% The two interrupt lines and the window they define.
\begin{tikzpicture}[font=\sffamily\scriptsize]

% ---- INT1, the watermark line: regular, frequent, tolerant of latency
\node[signame] at (-2mm,46mm) {INT1 watermark};
\draw[signal] (0mm,46mm) -- (120mm,46mm);
\foreach \x in {6,21,36,51,66,81,96,111}{
  \draw[signal] (\x mm,46mm) -- ++(0,4mm) -- ++(2mm,0) -- ++(0,-4mm);}

% ---- INT2, the high-g line: rare, and the moment everything is timed from
\node[signame] at (-2mm,34mm) {INT2 high-g};
\draw[signal] (0mm,34mm) -- (120mm,34mm);
\draw[signalb] (60mm,34mm) -- ++(0,4mm) -- ++(2mm,0) -- ++(0,-4mm);
\node[font=\sffamily\tiny, anchor=south] at (61mm,38.5mm) {threshold crossed in the sensor};

% ---- the window
\draw[tspan] (48mm,26mm) -- (60mm,26mm);
\node[tlabel] at (54mm,28mm) {pre 100 ms};
\draw[tspan] (60mm,26mm) -- (90mm,26mm);
\node[tlabel] at (75mm,28mm) {post 400 ms};

% ---- what the ring and the recorder are doing
\node[signame] at (-2mm,16.5mm) {ring};
\slice{14}{0}{60}{taskready}{fine channel fills the ring continuously}
\slice{14}{60}{30}{taskrun}{window held open}
\slice{14}{90}{30}{taskready}{ring resumes}

\node[signame] at (-2mm,4.5mm) {recorder};
\slice{2}{60}{14}{taskblock}{blanking}
\slice{2}{74}{16}{taskisr}{assemble}
\slice{2}{90}{22}{taskrun}{framed to USART3}

% ---- time axis and the one tick that matters
\draw[taxis] (0mm,-4mm) -- (124mm,-4mm) node[right, font=\tiny]{time};
\draw[tick] (60mm,-4mm) -- (60mm,52mm);
\node[font=\tiny, anchor=north] at (60mm,-4.5mm) {event tick};
\node[font=\tiny, anchor=north] at (0mm,-4.5mm) {0};

\node[font=\sffamily\tiny, text width=48mm, anchor=north west] at (0mm,-9mm)
  {Both marker pins are pulsed by their handlers and recorded on the
   instrument's digital inputs, time-aligned with its current trace. That is
   this bench's logic recorder; there is no oscilloscope and no logic analyser.};
\node[font=\sffamily\tiny, text width=48mm, anchor=north west] at (62mm,-9mm)
  {Every interval drawn here is a budget. None of them is measured until the
   cycle counter and the instrument have produced a number and the README names
   which one produced it.};
\end{tikzpicture}
